\section{Proposed Method}
\label{sec:method}

\subsection{AI Systems Being Tested}

This study will test several AI systems using two groups. The first group is made up of local models that run on a regular laptop with no internet connection, served through a free, open-source tool called Ollama \cite{ollama_2024}. Ollama makes it easy to download and run language models locally by handling all the setup automatically. The second group is made up of paid online AI services accessed through their APIs. Testing three different local models (instead of just one) is important because the results will then say something useful about local AI tools in general, not just about one specific model \cite{gemma4_phi4_qwen3_2026}.

\subsubsection*{Local Models (run via Ollama \cite{ollama_2024})}

\begin{enumerate}
    \item \textbf{\texttt{deepseek-ocr}} (DeepSeek AI) --- a document-focused vision-language model available on Ollama's model library, designed for reading and extracting text from images and scanned documents \cite{deepseek_ocr_2024}.
    \item \textbf{\texttt{gemma4:e4b}} (Google DeepMind, 2026) --- a local model with 4.5 billion active parameters that can read document images without a separate OCR tool and runs fully offline \cite{gemma4_2026}.
    \item \textbf{\texttt{qwen3.5:4b}} (Alibaba, 2025) --- a 4-billion-parameter model that is strong at understanding and pulling out structured information from documents \cite{qwen3_2025}.
    \item \textbf{\texttt{Phi-4-multimodal-instruct}} (Microsoft, 2025, \textit{optional}) --- a multimodal model that handles document images well through a multi-crop processing strategy \cite{phi4mini_2025}. This model will be included in the study if Ollama adds support for it; otherwise it will be evaluated as a secondary experiment.
\end{enumerate}

\subsubsection*{Cloud API Models}

\begin{enumerate}
    \item \textbf{GPT-4o} (OpenAI) --- a leading multimodal cloud model widely used for document understanding tasks, accessed via the OpenAI API \cite{openai_gpt4o_2024}.
    \item \textbf{Gemini 3.1 Pro} (Google DeepMind) --- Google's most capable multimodal cloud model, accessed via the Google AI API \cite{google_gemini_2025}.
    \item \textbf{Gemini 2.0 Flash} (Google DeepMind, \textit{optional third}) --- a faster and cheaper version of Gemini, included to show how a budget cloud option compares to the local models \cite{google_gemini_2025}.
\end{enumerate}

\subsection{Procedure}

Before any testing begins, 500 test documents will be created. These will look like real invoices and administrative forms but will use made-up data, so there are no privacy issues. For every document, a complete answer key will be prepared in advance that lists the exact correct answer for every field.

Each document will then be sent to every AI system three times using three different versions of the instruction. This is done because research has shown that AI models can give different results depending on how the instruction is written, so testing all three versions makes the comparison fair for every system \cite{anlsstar_2024, gemma4_phi4_qwen3_2026}:

\begin{enumerate}
    \item \textbf{Short instruction:} A single, simple sentence, for example: \emph{``Extract the vendor name, total amount, and date from this document.''}
    \item \textbf{Detailed instruction:} A longer description that explains every field, says what to do if a field is missing, and asks for a precise answer.
    \item \textbf{JSON instruction:} An instruction that asks the AI to return the answer as a structured JSON object only, for example: \emph{``Return only a JSON object with keys \texttt{vendor}, \texttt{total}, and \texttt{date}. Do not add any other text.''}
\end{enumerate}

Results will be recorded separately for all three instruction versions. For the main comparison between systems, the best-performing version for each system will be used. This way, no system is unfairly penalised because its instruction style happened to be a bad fit.

All local models will be downloaded and run through Ollama \cite{ollama_2024} on the same laptop, so speed and energy readings are directly comparable across models. The cloud models will be accessed through their respective APIs, and the exact token counts and costs will be logged from the billing records. A Python script will automatically check each AI's answers against the answer key, record whether each field was correct, and measure how long the AI took to respond.

\subsection{Measures}

\begin{itemize}
    \item \textbf{Accuracy (ANLS*):} Instead of a simple yes/no correct score, this study will use a metric called ANLS* \cite{anlsstar_2024}. This metric gives partial credit for answers that are almost correct --- for example, if the AI returns ``\$150.00'' when the answer is ``150.00'', it is not counted as a total failure. ANLS* also checks structured JSON output field by field, which gives a fairer and more useful picture of real performance. Results will be shown separately for each instruction version (short, detailed, JSON) and then summarised using the best version per system.

    \item \textbf{Speed:} How long each AI takes to process one document, measured from the moment the document is sent to the moment the final answer comes back. This will be recorded in milliseconds and averaged across all 500 documents for each system.

    \item \textbf{Cost Per Document:}
    \begin{itemize}
        \item \emph{Online AI:} The cost comes directly from the provider's billing records --- the number of tokens used, multiplied by the price per million tokens.
        \item \emph{Local AI models:} The cost has two parts:
        \begin{itemize}
            \item \textbf{Electricity:} The amount of power used during each test run will be measured using \texttt{powermetrics} on a Mac (this tool records CPU, GPU, and chip wattage in real time at 100\,ms intervals) or \texttt{powertop} on Linux \cite{tokenpowerbench_2024}. The power reading is then multiplied by the time taken and the local electricity rate (PKR/kWh) to get the energy cost per document.
            \item \textbf{Hardware share:} The laptop's purchase price is spread out evenly over three years (this is called straight-line depreciation) \cite{lenovo_tco_2026}. The per-document share is worked out based on how many documents are processed each day. This makes the cost calculation clear, repeatable, and easy for others to check.
        \end{itemize}
    \end{itemize}

    \item \textbf{Breakeven Point:} Using the cost numbers above, this is the number of documents per day at which the total yearly cost of a local model becomes equal to, and then lower than, the total yearly cost of the online service. Results will be shown as a simple graph for each local model compared to the online service.
\end{itemize}

\subsection{Test Documents}

All 500 test documents will be made from scratch to look like real invoices and forms, but using completely made-up information. This means there are no privacy concerns, and the right answers are already known before the test starts. The documents will have different layouts, different numbers of fields, and different types of content so that the results show how each AI performs across a range of realistic situations, not just one simple type of document \cite{bujang_samplesize_2021}.
